\documentclass[11pt]{article}
\usepackage[a4paper,margin=27mm]{geometry}
\usepackage{amsmath,amssymb,amsthm,xcolor,booktabs,array}
\newcommand{\PROVED}{\textcolor{green!40!black}{\textbf{PROVED}}}
\newcommand{\OPEN}{\textcolor{red!70!black}{\textbf{OPEN}}}
\newcommand{\AUDIT}{\textcolor{orange!70!black}{\textbf{AUDIT}}}
\newtheorem{theorem}{Theorem}
\newtheorem{corollary}{Corollary}
\title{Exact physical synthesis of the threshold delay channels (v111)}
\date{13 August 2026}

\begin{document}
\maketitle

\section{Arithmetic separation of a newborn cell}

Put
\begin{equation}
 T_N=\frac12\log N,
 \qquad
 \delta_N=T_{N+1}-T_N=\frac12\log\frac{N+1}{N},
 \qquad E_N=L^2(0,\delta_N).                                  \tag{1}
\end{equation}
For $1\le m\le N$ define two translated copies of the newborn boundary
cell:
\begin{align}
 L_m&=(\log m-T_{N+1},\,\log m-T_N),\notag\\
 R_m&=(T_N-\log m,\,T_{N+1}-\log m).                          \tag{2}
\end{align}
Let $J_m^L,J_m^R:E_N\to L^2(\mathbb R)$ be the translation isometries onto
$L^2(L_m)$ and $L^2(R_m)$, respectively.

\begin{theorem}[No-overlap theorem for an arithmetic threshold cell ---
\PROVED]
The $2N$ open intervals in (2) are pairwise disjoint.  Consequently the
delay synthesis
\begin{equation}
 \mathfrak S_N:
 \bigl(\ell^2\{1,\ldots,N\}\oplus\ell^2\{1,\ldots,N\}\bigr)
       \widehat\otimes E_N\longrightarrow L^2(\mathbb R),
\quad
 \mathfrak S_N(a,b;e)=
 \sum_{m\le N}a_mJ_m^Le+\sum_{m\le N}b_mJ_m^Re               \tag{3}
\end{equation}
is an isometry on elementary tensors and extends to an isometry on the
Hilbert tensor product.
\end{theorem}

\begin{proof}
All intervals have length $\delta_N$.  For consecutive integers
$m<m+1\le N$,
\[
 \log\frac{m+1}{m}\ge\log\frac N{N-1}
 >\log\frac{N+1}{N}=2\delta_N,
\]
so the intervals within each of the two families are disjoint.

An interval $L_m$ and an interval $R_n$ overlap in positive measure only
if both
\[
 \log m-T_{N+1}<T_{N+1}-\log n,
 \qquad
 T_N-\log n<\log m-T_N
\]
hold.  Using $2T_N=\log N$ and $2T_{N+1}=\log(N+1)$, these inequalities
are exactly
\[
                             N<mn<N+1,
\]
which is impossible for integers $m,n$.  Endpoint contacts have measure
zero.  Orthogonality of all ranges proves (3).
\end{proof}

Compression of (3) to any physical subwindow is a contraction.  Thus the
arithmetic label-to-delay map requested after the prime Julia colligation
is not an additional norm estimate: on every open threshold cell it is the
restriction of the explicit isometry $\mathfrak S_N$.

\section{The physical Witt word current}

For $k\ge1$ retain
\begin{equation}
 v_{N,k}(m)=\frac{\Lambda_k(m)}{\sqrt m(\log N)^k}.             \tag{4}
\end{equation}
Give different Witt depths orthogonal output copies and define
\begin{equation}
 \mathfrak W_Ne=
 \bigoplus_{k\ge1}\frac1{\sqrt2}
 \mathfrak S_N(v_{N,k},v_{N,k};e)
 \quad
 (e\in E_N).                                                   \tag{5}
\end{equation}

\begin{corollary}[Strict physical arithmetic current --- \PROVED]
The current (5) is constructed directly from the physical delay channels
and satisfies
\begin{equation}
 \|\mathfrak W_Ne\|^2
 =\left(\sum_{k\ge1}\frac{V_{N,k}}{(\log N)^{2k}}\right)\|e\|^2
 \longrightarrow
 \frac{\sqrt\pi}{2}e^{1/4}\operatorname{erf}\!\left(\frac12\right)
 \|e\|^2.                                                     \tag{6}
\end{equation}
Hence for some $N_0$ and $\eta>0$,
\begin{equation}
                         \|\mathfrak W_N\|^2\le1-\eta
 \qquad(N\ge N_0).                                           \tag{7}
\end{equation}
\end{corollary}

\begin{proof}
The two orientations in (5) have orthogonal physical ranges by Theorem 1,
and distinct depths were placed in orthogonal innovation channels.  Thus
the norm is the sum of the squared coefficient norms, giving the first
identity.  The fixed-depth and uniform all-depth arguments of v109 give
the limit in (6).  Its value is $0.5922965364\ldots<1$, proving (7).
\end{proof}

The boundary cross $H_{N,k}$ in v109 is now geometrically located: it is
an endpoint contact $mn=N$ or $mn=N+1$.  It disappears from the open-cell
$L^2$ synthesis because endpoints have measure zero, while it survives in
the differentiated threshold port G39.  This reconciles, without a sign
assumption, the word Gram of v110 with the physical delay current (5).

\section{Remaining Gamma--Tate coupling}

Theorem 1 closes the arithmetic label-to-physical-delay part of the
interface.  To identify (5) with the normalized Schur return one still has
to prove that the Kato-transported Gamma coherent cell and its affine
cocycle attach to $\mathfrak W_N$ through the already computed two Tate
terminal maps by a contraction
\begin{equation}
 \mathfrak C_N:
 \operatorname{Ran}\mathfrak W_N\oplus\mathscr K_{\Gamma,N}
 \oplus\mathbb C^2_{\rm Tate}longrightarrow
 \overline{\operatorname{Ran}D_0^{1/2}}                       \tag{8}
\end{equation}
whose value is
\begin{equation}
 \mathfrak C_N(\mathfrak W_Ne,\,b_{\Gamma,N}e,\,B'_{T,N}e)
 =D_0^{\dagger/2}(D_0h_e-q_e).                                \tag{9}
\end{equation}
The arithmetic restriction of (8) is now explicitly contractive by
(3)--(7).  What remains in (8)--(9) is the affine Gamma--Tate attachment,
not the Euler delay synthesis and not the all-depth unit budget.

\begin{center}
\begin{tabular}{>{\raggedright\arraybackslash}p{.55\linewidth}
                >{\centering\arraybackslash}p{.15\linewidth}
                >{\raggedright\arraybackslash}p{.20\linewidth}}
\toprule
Statement & Status & Location \\
\midrule
Exact disjointness of all newborn arithmetic delay channels & \PROVED & Theorem 1 \\
Physical label-to-delay isometry & \PROVED & (3) \\
Strict all-depth physical arithmetic current & \PROVED & (5)--(7) \\
Gamma--Tate attachment (8)--(9) & \OPEN & remaining interface \\
G40, condition $G$, row (d) & \OPEN & follow after the attachment and finite audit \\
\bottomrule
\end{tabular}
\end{center}

\end{document}
